\corrigesChapitre{11}{Graphes et matrices d’adjacence}
\begin{corrige}{1}{Vocabulaire}
\begin{questions}
\item Ordre $5$.
\item $14$ dans un graphe non orienté sans boucle.
\item Non ; $K_5$ aurait $5\times4/2=10$ arêtes, donc celui-ci n’est pas complet.
\end{questions}
\end{corrige}
\begin{corrige}{2}{Graphe du cours}
\begin{questions}
\item $A,B,D$.
\item $1$, car seule l’arête $CD$ lui est incidente.
\item Oui : les sommets du triangle sont reliés entre eux et $D$ rejoint $C$.
\end{questions}
\end{corrige}
\begin{corrige}{3}{Matrice d’un triangle}
\begin{questions}
\item $\begin{pmatrix}0&1&1\\1&0&1\\1&1&0\end{pmatrix}$.
\item $2$, le degré de chaque sommet.
\item $3\times2/2=3$.
\end{questions}
\end{corrige}
\begin{corrige}{4}{Une longueur}
\begin{questions}
\item $5$, car il y a cinq déplacements entre six sommets listés.
\item Oui. La répétition des sommets est autorisée pour une chaîne.
\item Non : l’arête $AC$ est utilisée plusieurs fois et toutes les arêtes ne sont pas parcourues exactement une fois.
\end{questions}
\end{corrige}
\begin{corrige}{5}{Calculer $A^2$}
\begin{questions}
\item $1$, par $A,C,D$.
\item $2$, via $B$ ou $C$.
\item Un parcours de longueur $2$ relie $A$ à $D$ ; deux reviennent en $A$. Les retours immédiats sont comptés.
\end{questions}
\end{corrige}
\begin{corrige}{6}{Un graphe orienté}
\begin{questions}
\item $A\to B$, $B\to C$, $C\to A$.
\item $I_3$, car trois transitions reviennent au point de départ.
\item $1$, puisque $M^4=M$.
\end{questions}
\end{corrige}
\begin{corrige}{7}{Complet}
\begin{questions}
\item $6\times5/2=15$.
\item Chaque sommet a degré $5$.
\item Non : les six degrés sont impairs ; le critère n’est pas satisfait.
\end{questions}
\end{corrige}
\begin{corrige}{8}{Une chaîne eulérienne}
\begin{questions}
\item $C$ de degré $3$ et $D$ de degré $1$.
\item $D,C,A,B,C$.
\item Non : ses extrémités sont les deux sommets impairs, distincts.
\end{questions}
\end{corrige}
\begin{corrige}{9}{Interpréter une puissance}
\begin{questions}
\item Il existe $17$ parcours orientés de longueur $5$ allant du sommet $2$ au sommet $3$.
\item Non : les produits matriciels comptent aussi les parcours répétés.
\item $0$, car les sommets $2$ et $3$ sont différents et aucun parcours vide ne les relie.
\end{questions}
\end{corrige}
\begin{corrige}{10}{Connexité et information}
\begin{questions}
\item Aucune arête ne relie un sommet du premier groupe à un sommet du second.
\item Non : ces deux groupes sont séparés.
\item Non : tous les produits conservent les blocs ; aucun parcours ne peut franchir une absence totale de liaison.
\end{questions}
\end{corrige}
\begin{corrige}{11}{Démontrer le comptage}
\begin{questions}
\item Il sépare un parcours en ses $n$ premières étapes et sa dernière étape.
\item Pour chaque parcours arrivant en $k$, la présence d’une dernière arête vers $j$ détermine s’il se prolonge ; dans un graphe simple, le facteur vaut $0$ ou $1$.
\item Les parcours classés par leur avant-dernier sommet sont disjoints et couvrent tous les cas. Cette somme est la formule du produit matriciel.
\end{questions}
\end{corrige}
\begin{corrige}{12}{Poids et comptage}
\begin{questions}
\item Non : l’interprétation de comptage exige la matrice d’adjacence appropriée.
\item Il somme des produits de poids le long des parcours, pas des sommes de distances ni un nombre de parcours.
\item Préciser ce que chaque coefficient représente et quelle opération correspond au problème. Un même graphe peut nécessiter plusieurs matrices selon la question.
\end{questions}
\end{corrige}
